\documentclass[12pt,a4paper]{article}
\usepackage[utf8]{inputenc}
\usepackage[T1]{fontenc}
\usepackage[italian]{babel}
\usepackage{geometry}
\usepackage{amsmath, amssymb}
\usepackage{listings}
\usepackage{xcolor}
\usepackage{hyperref}
\usepackage{booktabs}
\usepackage{array}
\usepackage{tikz}
\usetikzlibrary{arrows.meta, positioning, shapes.geometric, fit, backgrounds, matrix}
\usepackage{float}
\usepackage{caption}
\usepackage{enumitem}
\usepackage{tcolorbox}
\tcbuselibrary{breakable, skins}

\geometry{margin=2.5cm}

% --- Colori e stile ASP ---
\definecolor{aspkw}{RGB}{0,0,180}
\definecolor{aspcomment}{RGB}{60,120,60}
\definecolor{aspbg}{RGB}{250,252,255}
\definecolor{aspaggreg}{RGB}{160,0,160}

\lstdefinelanguage{ASP}{
  keywords={not, count, sum, min, max, minimize, maximize, show},
  keywordstyle=\color{aspkw}\bfseries,
  comment=[l]{\%},
  commentstyle=\color{aspcomment}\itshape,
  basicstyle=\ttfamily\small,
  backgroundcolor=\color{aspbg},
  frame=single,
  framerule=0.5pt,
  rulecolor=\color{gray!60},
  breaklines=true,
  showstringspaces=false,
  tabsize=4,
  numbers=left,
  numberstyle=\tiny\color{gray},
  numbersep=6pt,
  moredelim=[s][\color{aspaggreg}]{\#}{\ },
}

\lstset{language=ASP}

\newtcolorbox{defbox}[1]{colback=blue!5,colframe=blue!60!black,fonttitle=\bfseries,title=#1,breakable}
\newtcolorbox{exbox}[1]{colback=green!5,colframe=green!50!black,fonttitle=\bfseries,title=#1,breakable}
\newtcolorbox{warnbox}[1]{colback=orange!10,colframe=orange!80!black,fonttitle=\bfseries,title=#1,breakable}
\newtcolorbox{rulebox}[1]{colback=purple!5,colframe=purple!60!black,fonttitle=\bfseries,title=#1,breakable}

\hypersetup{colorlinks=true,linkcolor=blue!70!black}

\newcommand{\HRule}{\rule{\linewidth}{0.5mm}}

\begin{document}
\raggedbottom

% =====================================================================
% FRONTESPIZIO
% =====================================================================
\begin{titlepage}
  \centering
  \vspace*{2cm}

  {\Large\textbf{Universit\`a degli Studi di Torino}\\[0.3em]
  Corso di Laurea Magistrale in Informatica}

  \vspace{1cm}
  \HRule\\[0.4em]
  {\LARGE\textbf{Intelligenza Artificiale}}\\[0.3em]
  {\large Progetto Clingo}\\[0.2em]
  \HRule

  \vspace{1.8cm}
  {\huge\textbf{Pianificazione di Torneo Pok\'emon}}\\[0.6em]
  {\Large con Answer Set Programming}\\[0.4em]
  {\large Sistema Svizzero -- Relazione Tecnica Dettagliata}

  \vspace{3cm}
  {\large\textbf{Autore:}}\\[0.5em]
  {\large Alessandro Livero}

  \vspace{1.5cm}
  {\large Anno Accademico 2025/2026}

  \vfill
  {\large \today}
\end{titlepage}

\tableofcontents
\newpage

% =====================================================================
\section{Introduzione e Contesto}
% =====================================================================

\subsection{Il Sistema Svizzero}

Il \textbf{sistema svizzero} è un formato di torneo usato in scacchi, Pokémon
competitivo, Magic: The Gathering e altri giochi. A ogni turno i giocatori
vengono abbinati secondo queste regole:

\begin{enumerate}[noitemsep]
  \item Si abbinano giocatori con punteggio simile (\emph{spareggio}).
  \item Due giocatori non si affrontano mai due volte.
  \item Se il numero di giocatori è dispari, uno riceve il \textbf{BYE}
        (vittoria automatica senza giocare); nessuno lo può ricevere due volte.
\end{enumerate}

\subsection{Answer Set Programming (ASP)}

L'ASP è un paradigma di programmazione logica orientato alla soluzione di
problemi combinatori. Il \textbf{grounding + solving} di Clingo funziona in due fasi:

\begin{center}
\begin{tikzpicture}[
  box/.style={draw,rounded corners,minimum width=3cm,minimum height=1cm,align=center,font=\small},
  arrow/.style={-{Stealth},thick}
]
  \node[box,fill=blue!15] (ground) {Grounding\\\shortstack{\footnotesize sostituisce variabili\\con costanti}};
  \node[box,fill=green!15, right=2.5cm of ground] (solve) {Solving\\\shortstack{\footnotesize cerca \emph{answer sets}\\(modelli stabili)}};
  \node[box,fill=orange!15, right=2.5cm of solve] (opt) {Ottimizzazione\\\shortstack{\footnotesize seleziona i modelli\\con costo minimo}};
  \draw[arrow] (ground) -- (solve);
  \draw[arrow] (solve) -- (opt);
\end{tikzpicture}
\end{center}

Un \textbf{answer set} è un insieme di atomi ground che soddisfano tutte le
regole e i vincoli del programma. Con più answer set, Clingo può restituire
quello con il valore di ottimizzazione minimo (\texttt{--opt-mode=optN}).

\subsection{File del Progetto}

\begin{center}
\begin{tabular}{lll}
\toprule
\textbf{File} & \textbf{Tipo} & \textbf{Contenuto} \\
\midrule
\texttt{torneo.lp} & Programma ASP & Logica generica del torneo (invariante) \\
\texttt{turno1.lp} & Istanza & 5 giocatori, tutti a 0 punti, turno 1 \\
\texttt{turno2.lp} & Istanza & 8 giocatori, punteggi aggiornati, turno 2 \\
\bottomrule
\end{tabular}
\end{center}

\textbf{Uso da terminale:}
\begin{lstlisting}[numbers=none,backgroundcolor=\color{black!5}]
clingo torneo.lp turno1.lp 1 --opt-mode=optN
clingo torneo.lp turno2.lp 1 --opt-mode=optN
\end{lstlisting}
Il parametro \texttt{1} chiede la prima soluzione ottimale; \texttt{0} restituirebbe
tutte le soluzioni ottimali.

% =====================================================================
\section{Programma \texttt{torneo.lp} -- Analisi Dettagliata}
% =====================================================================

\subsection{Sezione 1: Normalizzazione e Ordinamento}

\subsubsection{\texttt{played\_already/2} -- Simmetria degli Incontri}

\begin{lstlisting}
played_already(P1, P2) :- played(P1, P2).
played_already(P1, P2) :- played(P2, P1).
\end{lstlisting}

Il fatto \texttt{played(p1, p2)} nell'istanza è \textbf{direzionale} (per
comodità di scrittura), ma semanticamente un incontro è simmetrico: se
Lorenzo ha giocato contro Emma, Emma ha giocato contro Lorenzo. \\
\texttt{played\_already/2} chiude questa relazione rispetto alla simmetria:
entrambi gli ordini sono derivati.

\subsubsection{\texttt{pid/2} -- Indice Canonico}

\begin{lstlisting}
pid(P, I) :- player(P),
             I = #count{ Q : player(Q), Q != P, Q < P }.
\end{lstlisting}

\begin{defbox}{Aggregato \texttt{\#count}}
\texttt{\#count\{ Q : Cond \}} conta il numero di valori distinti di \texttt{Q}
che soddisfano \texttt{Cond}. Il risultato è un intero $\ge 0$.
\end{defbox}

\textbf{Cosa calcola \texttt{pid}:}
\[
  \texttt{pid}(P, I) \iff I = |\{Q \in \text{player} \mid Q \neq P \land Q < P\}|
\]
Poiché Clingo ordina i simboli lessicograficamente (come stringhe), \texttt{pid}
assegna a ogni giocatore un indice intero univoco basato sull'ordine
lessicografico dei nomi. Con i giocatori \{diego, emma, giada, lorenzo, simone\}:

\begin{center}
\begin{tabular}{lc}
\toprule
\textbf{Giocatore} & \textbf{pid} \\
\midrule
diego & 0 \\
emma & 1 \\
giada & 2 \\
lorenzo & 3 \\
simone & 4 \\
\bottomrule
\end{tabular}
\end{center}

\textbf{Perché serve?} Per garantire che ogni coppia (P1, P2) sia considerata
\emph{una sola volta} nella generazione: solo la forma con
$\texttt{pid}(P1) < \texttt{pid}(P2)$ viene usata. Senza questo ordinamento,
Clingo genererebbe sia \texttt{match(emma, diego)} sia \texttt{match(diego, emma)}.

\subsection{Sezione 2: Generazione (Choice Rules)}

\subsubsection{\texttt{match/2} -- Generazione degli Incontri}

\begin{lstlisting}
{ match(P1, P2) } :-
    player(P1), player(P2),
    pid(P1, I1), pid(P2, I2), I1 < I2,
    not played_already(P1, P2).
\end{lstlisting}

\begin{defbox}{Choice Rule \texttt{\{ h \} :- Body}}
Una \textbf{choice rule} (regola di scelta) in ASP indica che l'atomo \texttt{h}
\emph{può} essere incluso nell'answer set se il corpo è soddisfatto,
ma non è obbligatorio. Clingo esplora entrambe le possibilità (incluso / non incluso)
e mantiene solo i modelli che soddisfano tutti i vincoli.
\end{defbox}

\textbf{Condizioni per la generazione di \texttt{match(P1, P2)}:}
\begin{enumerate}[noitemsep]
  \item Entrambi sono giocatori registrati.
  \item $\texttt{pid}(P1) < \texttt{pid}(P2)$: forma canonica (evita duplicati).
  \item \texttt{not played\_already(P1, P2)}: non si sono già affrontati.
\end{enumerate}

\subsubsection{\texttt{bye/1} -- Generazione del BYE}

\begin{lstlisting}
{ bye(P) } :- player(P), not had_bye(P).
\end{lstlisting}

Ogni giocatore che non ha mai ricevuto il BYE \emph{può} riceverlo in questo
turno. I vincoli hard successivi limitano a zero o uno i BYE effettivi.

\subsection{Sezione 3: Vincoli Hard (Integrity Constraints)}

\begin{defbox}{Integrity Constraint \texttt{:- Body}}
Un vincolo di integrità \texttt{:- Body} elimina dall'insieme degli answer set
qualsiasi modello in cui \texttt{Body} è vero. In sostanza: ``è
\textbf{vietato} che Body sia soddisfatto''.
\end{defbox}

\subsubsection{Helper \texttt{in\_match/1}}

\begin{lstlisting}
in_match(P) :- match(P, _).
in_match(P) :- match(_, P).
\end{lstlisting}

Deriva che \texttt{P} è impegnato in un match, indipendentemente dal ruolo
(primo o secondo argomento). Questo semplifica i vincoli successivi.

\subsubsection{Vincolo: Ogni Giocatore Deve Fare Qualcosa}

\begin{lstlisting}
:- player(P), not in_match(P), not bye(P).
:- player(P), in_match(P), bye(P).
\end{lstlisting}

\textbf{Prima riga}: è vietato che un giocatore non sia né in un match né abbia
il BYE. Ogni giocatore deve partecipare.

\textbf{Seconda riga}: è vietato che un giocatore sia contemporaneamente in un
match \emph{e} abbia il BYE. Le due alternative sono mutualmente esclusive.

\subsubsection{Vincolo: Un Giocatore in un Solo Match}

\begin{lstlisting}
:- match(P, Q1), match(P, Q2), Q1 != Q2.
:- match(P1, Q), match(P2, Q), P1 != P2.
:- match(P, _),  match(_, P).
\end{lstlisting}

Tre vincoli che coprono tutti i casi di conflitto:

\begin{center}
\renewcommand{\arraystretch}{1.3}
\begin{tabular}{cp{6.5cm}p{6cm}}
\toprule
\textbf{Vincolo} & \textbf{Proibisce} & \textbf{Esempio vietato} \\
\midrule
1 & P appare come P1 in due match diversi & \texttt{match(emma,diego)}, \texttt{match(emma,giada)} \\
2 & Q appare come P2 in due match diversi & \texttt{match(diego,emma)}, \texttt{match(giada,emma)} \\
3 & P appare in entrambi i ruoli & \texttt{match(emma,diego)}, \texttt{match(giada,emma)} \\
\bottomrule
\end{tabular}
\end{center}

\subsubsection{Vincolo: Un Solo BYE per Turno}

\begin{lstlisting}
:- 2 <= #count{ P : bye(P) }.
\end{lstlisting}

Vietato che due o più giocatori abbiano il BYE nello stesso turno. La sintassi
\texttt{N \textless= Agg} è equivalente a \texttt{Agg \textgreater= N}.

\subsubsection{Vincolo di Parità: BYE iff Numero Dispari}

\begin{lstlisting}
player_count(N) :- N = #count{ P : player(P) }.
:- player_count(N), 0 == N \ 2, bye(_).
:- player_count(N), 1 == N \ 2, not bye(_).
\end{lstlisting}

\textbf{\texttt{player\_count/1}}: deriva il numero totale di giocatori usando
\texttt{\#count}.

\textbf{\texttt{N \textbackslash{} 2}}: operatore modulo in Clingo (resto della
divisione intera).

\begin{itemize}[noitemsep]
  \item Se $N \mod 2 = 0$ (numero pari): vietato che esista un BYE.
  \item Se $N \mod 2 = 1$ (numero dispari): vietato che non esista un BYE.
\end{itemize}

Il BYE è obbligatorio esattamente quando il numero di giocatori è dispari.

\subsection{Sezione 4: Ottimizzazione}

\subsubsection{\texttt{score\_diff/3} -- Differenza di Punteggio}

\begin{lstlisting}
score_diff(P1, P2, D) :-
    match(P1, P2), score(P1, S1), score(P2, S2),
    S1 >= S2, D = S1 - S2.

score_diff(P1, P2, D) :-
    match(P1, P2), score(P1, S1), score(P2, S2),
    S2 > S1, D = S2 - S1.
\end{lstlisting}

Calcola la \textbf{differenza assoluta} di punteggio tra i due giocatori di ogni
match. Le due clausole coprono i due casi ($S1 \ge S2$ e $S2 > S1$) per evitare
valori negativi. Si noti che in ASP le espressioni aritmetiche come
\texttt{D = S1 - S2} sono valutate solo se \texttt{S1} e \texttt{S2} sono
già ground (istanziate).

\subsubsection{\texttt{\#minimize} -- Funzione di Ottimizzazione}

\begin{lstlisting}
#minimize { D, P1, P2 : score_diff(P1, P2, D) }.
\end{lstlisting}

\begin{defbox}{\texttt{\#minimize \{ peso, chiave1, chiave2, ... : Cond \}}}
Minimizza la \textbf{somma} dei pesi di tutti gli elementi dell'insieme che
soddisfano la condizione. La tupla \texttt{(peso, chiave1, chiave2, ...)}
identifica univocamente ogni elemento (le chiavi evitano conteggi doppi per
elementi con lo stesso peso ma identità diverse).
\end{defbox}

In questo caso: si minimizza $\sum_{(P1,P2) \in \text{match}} |S_{P1} - S_{P2}|$.

\textbf{Effetto pratico}: abbinamenti con giocatori a pari punteggio
hanno differenza 0 (contributo 0 alla somma) e sono preferiti; abbinamenti
tra giocatori a punteggio molto diverso sono penalizzati.

Questo implementa il principio base del sistema svizzero: gli abbinamenti
dello stesso ``scoreboard'' giocano tra loro.

\subsection{Sezione 5: Output}

\begin{lstlisting}
#show match/2.
#show bye/1.
\end{lstlisting}

\begin{defbox}{\texttt{\#show predicato/arità}}
Per default, Clingo mostra \emph{tutti} gli atomi dell'answer set. La direttiva
\texttt{\#show} limita l'output ai soli predicati elencati. Senza \texttt{\#show},
verrebbero visualizzati anche \texttt{pid}, \texttt{in\_match},
\texttt{played\_already}, ecc.
\end{defbox}

% =====================================================================
\section{Istanze: \texttt{turno1.lp} e \texttt{turno2.lp}}
% =====================================================================

\subsection{Turno 1 -- 5 Giocatori, Tutti a 0 Punti}

\begin{lstlisting}
player(lorenzo). player(emma). player(simone).
player(diego).   player(giada).

score(lorenzo, 0). score(emma, 0). score(simone, 0).
score(diego, 0).   score(giada, 0).
% Nessun played/2, nessun had_bye/1
\end{lstlisting}

\textbf{Situazione}: torneo appena iniziato.
\begin{itemize}[noitemsep]
  \item 5 giocatori (dispari): un BYE obbligatorio.
  \item Tutti a 0 punti: la differenza di punteggio è 0 per qualsiasi abbinamento.
  \item Nessun incontro già avvenuto: tutti gli abbinamenti sono validi.
  \item Nessun BYE precedente: chiunque può riceverlo.
\end{itemize}

Con tutti i punteggi uguali, l'ottimizzatore minimizza una somma di zeri:
\emph{qualsiasi} abbinamento è ottimo. Clingo restituirà una soluzione arbitraria
tra quelle valide (dipende dall'ordine di esplorazione).

\begin{exbox}{Output reale -- \texttt{clingo torneo.lp turno1.lp 1 --opt-mode=optN}}
\begin{verbatim}
Answer: 1
bye(simone) match(giada,lorenzo) match(diego,emma)
Optimization: 0
OPTIMUM FOUND
\end{verbatim}
La funzione obiettivo vale \textbf{0}, come previsto: tutti i giocatori sono a
pari punteggio, quindi ogni possibile combinazione di 2 match + 1 bye è
equivalentemente ottima (Clingo ne trova più di una con lo stesso costo
durante la prova di ottimalità).
\end{exbox}

\subsection{Turno 2 -- 8 Giocatori, Punteggi Differenziati}

\begin{lstlisting}
% Vecchi giocatori con punteggi aggiornati
player(lorenzo). score(lorenzo, 2).
player(simone).  score(simone,  2).
player(giada).   score(giada,   2).
player(emma).    score(emma,    0).
player(diego).   score(diego,   0).
% Nuovi iscritti (tutti a 0)
player(jacopo).  score(jacopo,  0).
player(thomas).  score(thomas,  0).
player(desiree). score(desiree, 0).

% Incontri del turno 1
played(emma, lorenzo).
played(diego, simone).
% had_bye: giada non puo' ricevere di nuovo il BYE
had_bye(giada).
\end{lstlisting}

\textbf{Analisi}:
\begin{itemize}[noitemsep]
  \item 8 giocatori (pari): nessun BYE questo turno.
  \item I giocatori si raggruppano per punteggio: $\{2\}$ = \{lorenzo, simone, giada\}
        e $\{0\}$ = \{emma, diego, jacopo, thomas, desiree\}.
  \item \texttt{played(emma, lorenzo)} e \texttt{played(diego, simone)}: questi
        abbinamenti non possono ripetersi.
  \item L'ottimizzatore cercherà di abbinare i ``2'' tra loro e i ``0'' tra loro.
\end{itemize}

\begin{exbox}{Output reale -- \texttt{clingo torneo.lp turno2.lp 1 --opt-mode=optN}}
\begin{verbatim}
Answer: 1
match(giada,lorenzo)  match(desiree,simone)
match(diego,jacopo)   match(emma,thomas)
Optimization: 2
OPTIMUM FOUND
\end{verbatim}
Come previsto, \texttt{giada} e \texttt{lorenzo} (entrambi a 2 punti) vengono
abbinati tra loro a costo 0. Poiché i giocatori a 2 punti sono 3 (numero
dispari nel loro gruppo), \texttt{simone} resta senza un pari-punteggio
disponibile e viene abbinato a \texttt{desiree} (0 punti): questo
\emph{``across-group'' match} contribuisce $|2-0|=2$ alla somma, che è
esattamente il valore minimo raggiungibile. Gli altri giocatori a 0 punti
(\texttt{diego}, \texttt{emma}, \texttt{jacopo}, \texttt{thomas}) si abbinano
liberamente tra loro a costo 0, rispettando i vincoli su \texttt{played\_already}.
\end{exbox}

% =====================================================================
\section{Concetti ASP Fondamentali Usati nel Progetto}
% =====================================================================

\subsection{Riepilogo dei Costrutti}

\begin{center}
\renewcommand{\arraystretch}{1.4}
\begin{tabular}{p{4.5cm}p{4cm}p{5cm}}
\toprule
\textbf{Costrutto} & \textbf{Sintassi} & \textbf{Semantica} \\
\midrule
Fatto & \texttt{p(a).} & \texttt{p(a)} è sempre vero \\
Regola normale & \texttt{h :- b1, b2.} & \texttt{h} è vero se \texttt{b1} e \texttt{b2} sono veri \\
Negazione per fallimento & \texttt{not p} & vero se \texttt{p} non è nell'answer set \\
Choice rule & \texttt{\{h\} :- b.} & \texttt{h} può o non può essere incluso \\
Vincolo hard & \texttt{:- b.} & elimina i modelli dove \texttt{b} è vero \\
Aggregato count & \texttt{\#count\{x : p(x)\}} & numero di \texttt{x} con \texttt{p(x)} \\
Modulo & \texttt{N \textbackslash{} M} & resto di N diviso M \\
Minimizzazione & \texttt{\#minimize\{w,k : p\}.} & minimizza la somma dei pesi \\
Direttiva show & \texttt{\#show p/n.} & mostra solo il predicato \texttt{p} \\
\bottomrule
\end{tabular}
\end{center}

\subsection{Ciclo di Vita di un Answer Set}

\begin{enumerate}
  \item \textbf{Grounding}: ogni variabile è sostituita con tutte le costanti
        compatibili. Le regole con variabili diventano insiemi di regole ground.
  \item \textbf{Solving}: il solver esplora esponenzialmente le possibili
        combinazioni di atomi. Per ogni combinazione verifica:
        \begin{enumerate}[noitemsep]
          \item Tutte le regole normali sono soddisfatte (le teste sono derivate
                quando i corpi sono veri).
          \item Tutti i vincoli hard sono soddisfatti (nessun corpo vietato è
                vero).
          \item Il modello è \emph{minimo}: nessun sottoinsieme è ancora un
                modello stabile.
        \end{enumerate}
  \item \textbf{Ottimizzazione}: tra tutti gli answer set validi, Clingo
        seleziona quelli con il valore della funzione obiettivo minimo.
\end{enumerate}

\subsection{Negazione per Fallimento vs Negazione Classica}

\begin{warnbox}{Differenza critica}
\textbf{Negazione classica} (\texttt{-p}): afferma esplicitamente la falsità di
\texttt{p}. Richiede che il programma contenga \texttt{-p} come fatto o
conclusione.

\textbf{Negazione per fallimento} (\texttt{not p}): vera se \texttt{p}
\emph{non appare} nell'answer set corrente (assunzione del mondo chiuso).

Nel progetto si usa solo \texttt{not}, appropriato perché: se non è noto che
un giocatore abbia avuto il BYE (\texttt{had\_bye}), si assume che non lo
abbia avuto.
\end{warnbox}

% =====================================================================
\section{Architettura Complessiva e Flusso}
% =====================================================================

\begin{center}
\begin{tikzpicture}[
  node distance=1.2cm,
  box/.style={draw,rounded corners,align=center,minimum width=3.5cm,
              minimum height=0.9cm,font=\small},
  arrow/.style={-{Stealth},thick}
]
  \node[box,fill=gray!15] (input) {Input istanza\\{\footnotesize player, score, played, had\_bye}};
  \node[box,fill=blue!15, below=of input] (norm) {Normalizzazione\\{\footnotesize played\_already, pid}};
  \node[box,fill=green!15, below=of norm] (gen) {Generazione\\{\footnotesize \{match\}, \{bye\}}};
  \node[box,fill=red!15, below=of gen] (constr) {Vincoli Hard\\{\footnotesize partecipazione, unicità, parità}};
  \node[box,fill=purple!15, below=of constr] (opt) {Ottimizzazione\\{\footnotesize \#minimize score\_diff}};
  \node[box,fill=orange!15, below=of opt] (out) {Output\\{\footnotesize match/2, bye/1}};

  \draw[arrow] (input) -- (norm);
  \draw[arrow] (norm) -- (gen);
  \draw[arrow] (gen) -- (constr);
  \draw[arrow] (constr) -- (opt);
  \draw[arrow] (opt) -- (out);

  \node[right=1.5cm of constr, font=\footnotesize, text=gray, align=left]
    (note) {Elimina modelli\\non validi};
  \draw[->,gray,dashed] (note) -- (constr);
\end{tikzpicture}
\end{center}

\subsection{Separazione Programma / Istanza}

La separazione tra \texttt{torneo.lp} (logica) e i file di istanza
(\texttt{turno1.lp}, \texttt{turno2.lp}) è il punto di forza del paradigma ASP:
\begin{itemize}[noitemsep]
  \item Per aggiungere un nuovo turno basta creare un nuovo file istanza.
  \item La logica del torneo non va mai modificata.
  \item Diversi tornei (con regole diverse) potrebbero condividere le istanze.
\end{itemize}

% =====================================================================
\section{Proprietà Formali e Correttezza}
% =====================================================================

\subsection{Completezza: Ogni Abbinamento Valido è Generato}

La choice rule \texttt{\{match(P1,P2)\}} genera \emph{ogni} coppia di giocatori
che non si siano già affrontati (nella forma canonica). I vincoli hard poi
filtrano le combinazioni non valide. Il solver esplora tutte le possibilità:
nessun abbinamento valido viene perduto.

\subsection{Unicità della Forma Canonica}

L'uso di \texttt{pid(P1, I1), pid(P2, I2), I1 < I2} garantisce che ogni coppia
$\{P1, P2\}$ appaia esattamente una volta nella generazione. Senza questo
ordinamento, il grounding produrrebbe sia \texttt{match(A,B)} che
\texttt{match(B,A)} come atomi indipendenti, e i vincoli dovrebbero essere
molto più complessi per eliminarli.

\subsection{Ottimalità dell'Abbinamento}

Il \texttt{\#minimize} non garantisce che l'abbinamento sia ``perfetto'' nel
senso matematico (problema di accoppiamento pesato su grafo, NP-completo in
generale), ma trova il \textbf{minimo globale} della somma delle differenze,
che è la definizione standard di abbinamento ottimale nel sistema svizzero.

\end{document}
